%-------------------------
% Resume in LaTeX (Korean)
% Based on: https://github.com/sb2nov/resume
% License: MIT
%------------------------

\documentclass[letterpaper,11pt]{article}

\usepackage[empty]{fullpage}
\usepackage{titlesec}
\usepackage[usenames,dvipsnames]{color}
\usepackage{enumitem}
\usepackage[hidelinks]{hyperref}
\usepackage{fancyhdr}
\usepackage{tabularx}
\usepackage{kotex}
\usepackage{iftex}
\ifpdftex
  \input{glyphtounicode}
\fi

\pagestyle{fancy}
\fancyhf{}
\fancyfoot{}
\renewcommand{\headrulewidth}{0pt}
\renewcommand{\footrulewidth}{0pt}

\addtolength{\oddsidemargin}{-0.5in}
\addtolength{\evensidemargin}{-0.5in}
\addtolength{\textwidth}{1in}
\addtolength{\topmargin}{-.5in}
\addtolength{\textheight}{1.0in}

\urlstyle{same}
\raggedbottom
\raggedright
\setlength{\tabcolsep}{0in}

\titleformat{\section}{
  \vspace{-4pt}\scshape\raggedright\large
}{}{0em}{}[\color{black}\titlerule \vspace{-5pt}]

\ifpdftex
  \pdfgentounicode=1
\fi

%-------------------------
% Custom commands
\newcommand{\resumeItem}[2]{
  \item\small{
    \textbf{#1}{: #2 \vspace{-2pt}}
  }
}

\newcommand{\resumeSubheading}[4]{
  \vspace{-1pt}\item
    \begin{tabular*}{0.97\textwidth}[t]{l@{\extracolsep{\fill}}r}
      \textbf{#1} & #2 \\
      \textit{\small #3} & \textit{\small #4} \\
    \end{tabular*}\vspace{-5pt}
}

\newcommand{\resumeSubSubheading}[2]{
    \begin{tabular*}{0.97\textwidth}{l@{\extracolsep{\fill}}r}
      \textit{\small #1} & \textit{\small #2} \\
    \end{tabular*}\vspace{-5pt}
}

\newcommand{\resumeProjectSubheading}[2]{
    \vspace{2pt}
    \begin{tabular*}{0.97\textwidth}{l@{\extracolsep{\fill}}r}
      \small\textbf{#1} & \textit{\small #2} \\
    \end{tabular*}\vspace{-5pt}
}

\newcommand{\resumeProjectHeading}[2]{
    \item
    \begin{tabular*}{0.97\textwidth}{l@{\extracolsep{\fill}}r}
      \small #1 & #2 \\
    \end{tabular*}\vspace{-7pt}
}

\newcommand{\resumeItemPlain}[1]{
  \item\small{
    {#1 \vspace{-2pt}}
  }
}

\renewcommand{\labelitemii}{$\circ$}

\newcommand{\resumeSubHeadingListStart}{\begin{itemize}[leftmargin=*]}
\newcommand{\resumeSubHeadingListEnd}{\end{itemize}}
\newcommand{\resumeItemListStart}{\begin{itemize}}
\newcommand{\resumeItemListEnd}{\end{itemize}\vspace{-5pt}}

%-------------------------------------------
\begin{document}

%----------HEADING-----------------
\begin{tabular*}{\textwidth}{l@{\extracolsep{\fill}}r}
  \textbf{\Large 문수영 (Sooyoung Moon)} & 이메일: \href{mailto:msy0128@gmail.com}{msy0128@gmail.com} \\
  \href{https://github.com/symoon94}{github.com/symoon94} & 연락처: +82-10-9473-5586 \\
\end{tabular*}


%-----------EXPERIENCE-----------------
\section{경력}
  \resumeSubHeadingListStart

    \resumeSubheading
      {\href{https://www.upstage.ai/}{Upstage}}{원격 근무}
      {팀 리드, SolarBox}{2026.04 -- 현재}
    \resumeSubSubheading
      {테크 리드, SolarBox}{2025.08 -- 2026.03}
    \resumeSubSubheading
      {소프트웨어 엔지니어}{2021.06 -- 2025.07}

      \resumeProjectSubheading{온프레미스 AI 솔루션 플랫폼}{2024.12 -- 현재}
      \resumeItemListStart
        \resumeItem{자율 팀 에이전트 (Milchick)}
          {4개 이상의 팀에 걸친 수동 브리핑과 Q\&A 라우팅이 상당한 조율 오버헤드를 만들고 있었음. 비용과 품질의 균형을 위해 듀얼 LLM 아키텍처를 선택---\textbf{Solar}로 빠른 의도 분류, \textbf{Claude}로 복잡한 도구 실행. 품질 보증을 위한 eval 테스트 스위트, 선택적 컨텍스트 로딩이 가능한 지식 시스템, 내부 데이터 anti-leak 보호 기능을 구현. 계정 생성, 라이선스 발급, 카탈로그 관리를 위한 재사용 가능 에이전트 스킬도 개발. 현재 SolarBox 팀 전체가 매일 브리핑, Q\&A, 파트너 운영에 사용 중.}
        \resumeItem{엔터프라이즈 포탈 개발 리드}
          {파트너사가 AI 솔루션 패키지에 접근할 셀프서비스 수단이 없어 고객별 수동 전달이 필요했음. 요구사항 수집부터 프로덕션까지 개발을 주도: \textbf{AWS CDK} 개발/운영 환경 분리, \textbf{GitHub Actions} CI/CD, \textbf{i18n}(83개 파일 영문 번역), 역할 기반 draft 버전 관리, ``Edit with Claude'' AI 편집 기능. AWS, Azure, Snowflake에 걸친 글로벌 파트너 셀프서비스 실현.}
        \resumeItem{SolarBox 카탈로그 및 릴리스 관리}
          {멀티테넌트 환경에서 80개 이상의 AI 제품 카탈로그를 관리할 중앙화된 시스템이 부재했음. 검색, 페이지네이션, CUDA suffix 기반 시맨틱 버전 비교가 가능한 manifest API를 구축하고 \textbf{Keycloak} SSO와 \textbf{ClickHouse} 마이그레이션 트래킹을 설정. 4개 이상의 엔지니어링 팀 간 SolarBox 2.5 릴리스를 조율하여 6개 AI 제품을 3개 클라우드 마켓플레이스에 동시 출시.}
        \resumeItem{모델 암호화 파이프라인}
          {수동 암호화 단계가 카탈로그 등록의 병목이 되고 있었음. GitHub Workflow, \textbf{Kubernetes} Job, \textbf{Ansible}을 활용한 end-to-end LLM 모델 암호화를 자동화하고, \textbf{AWS OIDC}로 안전한 S3 아티팩트 관리 및 체크섬 검증을 구현. 암호화부터 업로드, 검증까지 모든 수동 개입을 제거.}
      \resumeItemListEnd

      \resumeProjectSubheading{Document AI 서빙 시스템}{2023.07 -- 현재}
      \resumeItemListStart
        \resumeItem{비동기 추론 시스템}
          {동기 추론만으로는 엔터프라이즈 규모의 문서를 처리할 수 없었음. v1은 \textbf{SQS}/Lambda(SaaS)와 RabbitMQ(온프렘)로 설계하고, 운영 복잡도를 낮추기 위해 v2에서 \textbf{Asynq}(Go 기반 Redis 큐)로 전환. 배포마다 CI 연동 sanity test를 실행하고 p1/p2/p3 알람 체계를 구축하여 실행 신뢰성 확보. 비동기 응답 시간 30초에서 0.2초로, 88페이지 문서 처리 시간 21분에서 4분으로 단축. 현재 일 160,000장 이상 처리하며 기존 대비 \textbf{18배 쓰루풋} 달성.}
        \resumeItem{서비스 안정성}
          {10개월간 팀 전체가 해결하지 못한 간헐적 패닉이 있었음. 체계적인 재현을 통해 근본 원인이 \textbf{go-fitz} PDF 라이브러리의 스레드 비안전 내부 구현임을 추적하고, 애플리케이션 레벨 \textbf{Locking}을 적용하여 해결. 온프렘 멀티노드 환경에 \textbf{Prometheus}/Grafana 모니터링을 구축. 일 10,000건 이상의 Segmentation Fault를 완전히 제거하고 500 에러 오알람을 주 65건에서 0건으로 감소.}
      \resumeItemListEnd

      \resumeProjectSubheading{Document AI MLOps 플랫폼}{2022.02 -- 2023.08}
      \resumeItemListStart
        \resumeItem{스토리지 최적화}
          {일 20,000건 이상의 이미지를 적재하는 고객 환경에서 UI 로드 시간이 허용 불가 수준이었음. 쿼리 패턴을 분석하고 이미지 인입 파이프라인에 타겟 인덱싱을 적용하여 로드 시간 \textbf{98\%} 감소 달성.}
        \resumeItem{실시간 이벤트 시스템}
          {모델 학습, 평가, 추론 작업의 진행 상황을 확인할 수 없어 사용자가 상태를 폴링해야 했음. \textbf{SSE} 기반 실시간 이벤트 아키텍처를 구현하여 라이브 추적 가능하게 함.}
      \resumeItemListEnd

  \resumeSubHeadingListEnd


%-----------EDUCATION-----------------
\section{학력}
  \resumeSubHeadingListStart
    \resumeSubheading
      {University of Arizona}{Tucson, AZ}
      {컴퓨터과학 학사}{2021.05}
    \resumeSubheading
      {국민대학교}{서울}
      {발효융합공학 학사}{2017.08}
  \resumeSubHeadingListEnd


%-----------PRESENTATIONS-----------------
\section{발표}
  \resumeSubHeadingListStart
    \resumeProjectHeading
      {\textbf{신입에서 리더까지: 주니어의 관성을 깨며 배운 것들} $|$ \emph{\href{https://www.slideshare.net/slideshow/from-junior-developer-to-tech-lead-leveraging-ai-for-leadership-and-team-impact/286718141}{Women Techmakers Korea 2026}}}{2026.03}
    \resumeProjectHeading
      {\textbf{AI로 굴러가게 만들기: 작은 팀의 생존 기술} $|$ \emph{\href{https://www.slideshare.net/slideshow/gdg-incheon-2025-ai/284715174}{GDG Incheon DevFest 2025}}}{2025.11}
    \resumeProjectHeading
      {\textbf{취업 준비부터 AI 스타트업 5년차까지} $|$ \emph{\href{https://www.slideshare.net/slideshow/2025-ai-5/284715261}{중앙대학교 커리어 세미나}}}{2025.11}
    \resumeProjectHeading
      {\textbf{Document AI 서비스 성장 스토리} $|$ \emph{\href{https://www.slideshare.net/slideshow/mlops-api-solarbox-document-ai/283755773}{제주 AI 컨퍼런스}} $|$ \emph{\href{https://www.hanbit.co.kr/product/C2486251278}{한빛+ DevGround 2025}}}{2025.10}
  \resumeSubHeadingListEnd


%-------------------------------------------
\end{document}
